%
% MODELO DE PLANO DE ENSINO — FACULDADE DE ENGENHARIA ELÉTRICA (FEE)
% Campus Universitário de Tucuruí — Universidade Federal do Pará (UFPA)
%
% Este modelo de Plano de Ensino foi preparado pelo docente
% André Felipe Souza da Cruz, no âmbito da Faculdade de Engenharia Elétrica,
% para padronização e apoio à elaboração dos Planos de Ensino das Atividades
% Curriculares, em conformidade com a legislação e as normas institucionais da UFPA.
%
% -------------------------------------------------------------------------
% CRÉDITOS
%
% Este template é uma adaptação do modelo originalmente disponibilizado por
% Rodrigo Cavalcante, disponível em:
% https://www.overleaf.com/articles/plano-de-aula/chhjxtdqjsqk
%
% A presente versão mantém os créditos ao autor original, incorporando
% adaptações estruturais, organizacionais e normativas voltadas ao uso
% institucional pela Faculdade de Engenharia Elétrica da UFPA.
%
% -------------------------------------------------------------------------
% ORGANIZAÇÃO DO PROJETO
%
% Este arquivo (PlanoEnsino.tex) concentra o cabeçalho, a identificação da
% disciplina, a tabela de informações gerais, o texto (objetivos,
% competências, conteúdo programático, metodologia, avaliação), a
% programação semestral e os anexos.
%
% A bibliografia deve ser inserida no arquivo referencias.bib, utilizando
% os comandos do pacote biblatex-abnt, adotado para adequação às normas da ABNT.
%
% Os pacotes utilizados neste modelo estão organizados no arquivo Package.tex.
%
% -------------------------------------------------------------------------
% OBSERVAÇÕES IMPORTANTES
% - Caso o documento seja utilizado como Plano de ENSINO (e não Plano de Aula),
% alterar o termo correspondente logo abaixo.
% - Para habilitar a seção "Programação Semestral", que contém a tabela com o
% detalhamento do conteúdo por aula, remover o símbolo "%" da linha que
% inclui essa seção, mais abaixo neste arquivo.
%
% -------------------------------------------------------------------------

\documentclass[oneside,a4paper,12pt]{article}
\usepackage[english,brazilian]{babel}
\usepackage[utf8]{inputenc}
\usepackage[T1]{fontenc}
%\usepackage{lmodern} %{lmodern}	% Usa a fonte Latin Modern			
\usepackage[top=20mm, bottom=20mm, left=20mm, right=20mm]{geometry}
\usepackage{framed}
\usepackage{booktabs} % Pacote para deixar tabelas mais bonitas.
\usepackage{color}	% Pacote de Cores
\usepackage[colorlinks=true,linkcolor=black,anchorcolor=black,citecolor=black,filecolor=black,menucolor=black,runcolor=black,urlcolor=black]{hyperref}%
\usepackage{graphicx}	% Pacote de imagens
\usepackage{microtype} 	% para melhorias de justificação
\usepackage{amsmath}
\usepackage[hang, small,labelfont=bf,up,textfont=it,up]{caption} % Custom captions under/above floats in tables or figures
\usepackage{indentfirst}% Indenta o primeiro parágrafo de cada seção.
\usepackage{enumitem} % Customized lists
\setlist[itemize]{noitemsep} % Make itemize lists more compact
\usepackage{array} %permite o uso de tabelas
\usepackage{longtable} %permite o uso de tabelas personalizadas
\usepackage{rotating} %permite rotacionar imagens
\usepackage{tikz} %permite o desenho de formas geométricas simples dadas suas características
\usepackage{ragged2e}%
\usepackage{csquotes}%
%\usepackage{setspace}%
\usepackage[style=abnt]{biblatex} %coloque backref=true se quiser que fale em qual pagina foi citado
\addbibresource{referencias.bib} % Seus arquivos de bibliografia
\definecolor{shadecolor}{rgb}{0.9,0.9,0.9}
\graphicspath{ {./Images/} }

% Definições-----------------------------------------------------------
\newcommand{\plano}{Ensino}%Ensino ou Aula
\newcommand{\universidade}{Universidade Federal do Pará \par Campus Universitário de Tucuruí }
\newcommand{\departamento}{Faculdade de Engenharia Elétrica}

% ===== Dados da Disciplina =====
\newcommand{\professor}{Nome do(a) docente responsável}
\newcommand{\email}{email@ufpa.br}
\newcommand{\curso}{Nome do Curso}
\newcommand{\disciplina}{Nome da Disciplina}
\newcommand{\carater}{Obrigatória ou Optativa}
\newcommand{\cargahoraria}{Carga horária total}
\newcommand{\regime}{Regime de oferta}
\newcommand{\periodo}{Período de oferta}
\newcommand{\prerequisito}{Pré-requisito(s), se houver}
\newcommand{\periodoletivo}{Ano.Período}

% ---------------------------------------------------------------------
\begin{document}
  \pagestyle{empty}
% Cabeçalho--------------------------------------------------------------
\begin{center}
    \includegraphics[width=2.5cm]{Figuras/UFPA.png}

    \vspace{8pt}

    \universidade
    \par
    \departamento
    \par

    \vspace{16pt}

    \LARGE \textbf{Plano de \plano}
\end{center}

%\vspace{6pt}
%TABELA----------------------------------------------------------------
\begin{center}
\renewcommand{\arraystretch}{1.2}

\begin{tabular}{|p{0.48\textwidth}|p{0.48\textwidth}|}
\hline
\multicolumn{2}{|c|}{\textbf{Dados de Identificação}} \\
\hline
\textbf{Professor:} Nome do(a) docente responsável &
\textbf{E-mail:} email@instituicao.br \\
\hline
\multicolumn{2}{|l|}{\textbf{Curso:} Nome do Curso} \\
\hline
\textbf{Disciplina:} Código - Nome da Disciplina &
\textbf{Caráter:} Obrigatória ou Optativa \\
\hline
\textbf{Carga horária:} XXh &
\textbf{Regime:} Regime de oferta \\
\hline
\multicolumn{2}{|l|}{\textbf{Período de oferta:} dd/mm/aaaa à dd/mm/aaaa} \\
\hline
\multicolumn{2}{|l|}{\textbf{Pré-requisito:} Informar pré-requisito(s), se houver} \\
\hline
\multicolumn{2}{|l|}{\textbf{Período letivo:} Ano.Período} \\
\hline
\end{tabular}
\end{center}

\vspace{3mm}
%Texto-----------------------------------------------------------------
  \begin{snugshade}
  \section{Ementa}
  \end{snugshade}

  Escreva aqui a ementa da disciplina conforme descrito no Projeto Pedagótico de Curso de Engenharia Elétrica, aprovado em 2012. Alterações da ementa devem ser comunicados à direção da FEE para apreciação do colegiado.

  \begin{snugshade}
  \section{Objetivos}
\end{snugshade}

O objetivo da disciplina é proporcionar ao discente a compreensão dos conceitos fundamentais da área abordada, bem como o desenvolvimento de habilidades para análise, interpretação e resolução de problemas teóricos e práticos. Especificamente, busca-se que o discente seja capaz de aplicar métodos e técnicas pertinentes à área, interpretar resultados em diferentes contextos, estabelecer relações entre teoria e prática e desenvolver raciocínio lógico, pensamento crítico e autonomia no processo de aprendizagem.

\begin{snugshade}
  \section{Competências e Habilidades}
\end{snugshade}

\begin{itemize}

    \item Compreender e aplicar os conceitos fundamentais da área abordada na disciplina. Relacionar competências a serem desenvolvidas conforme seção 3.4 Competências e Habilidades, página 13, do PPC de Engenharia Elétrica de 2012.

    \item Analisar e interpretar problemas teóricos e práticos, estabelecendo relações com aplicações na Engenharia.

    \item Utilizar métodos e técnicas adequados para modelagem e resolução de problemas.

    \item Desenvolver raciocínio lógico, pensamento crítico e rigor na análise de resultados.

\end{itemize}

Ao final da disciplina, espera-se que o discente seja capaz de aplicar os conceitos estudados na análise e resolução de problemas, interpretar resultados obtidos em diferentes contextos e utilizar ferramentas e métodos adequados à área. Espera-se, ainda, que o discente demonstre autonomia na aprendizagem, capacidade de modelagem e habilidade de relacionar os conteúdos teóricos com aplicações práticas em Engenharia.


  \begin{snugshade}
  \section{Conteúdo programático} % conteúdos programados para a aula organizados em tópicos (de 4 a 8).
  \end{snugshade}

  O conteúdo programático deve ser elaborado em conformidade com a ementa da disciplina prevista no PPC do curso, assegurando coerência entre os tópicos abordados, os objetivos de aprendizagem e a carga horária disponível.

  Recomenda-se que os conteúdos sejam organizados em unidades temáticas, apresentadas em sequência lógica e progressiva.

  Cada unidade deve possuir título e descrição sucinta dos conteúdos a serem desenvolvidos ao longo do período letivo.



  \begin{snugshade}
  \section{Procedimentos metodológicos} % estratégias relevantes adotadas para alcançar os objetivos.
  \end{snugshade}

A metodologia de ensino deve apresentar, de forma sucinta, as estratégias pedagógicas e os recursos didáticos utilizados no desenvolvimento da disciplina, tais como aulas expositivas e dialogadas, resolução de exercícios, seminários, estudos de caso e atividades individuais ou em grupo. Recomenda-se também indicar os ambientes de realização das atividades, como sala de aula, laboratório de informática ou laboratório prático, bem como o uso de recursos e plataformas digitais de apoio ao ensino.


  \begin{snugshade}
  \section{Recursos didáticos} % quadro, giz, retro-projetor, filme, música, quadrinhos, etc.
  \end{snugshade}
  Exemplos:
    \begin{itemize}
	  \item apoio de recursos audiovisuais (data-show e filmes).
	  \item laboratório de informática ou= laboratório multidisciplinar.
	  \item Google Sala de Aula (Classroom)
    \end{itemize}
  %\newpage
  \begin{snugshade}
\section{Avaliação}
\end{snugshade}

A avaliação da aprendizagem será realizada de forma contínua e processual, por meio de instrumentos diversificados, tais como provas escritas, trabalhos individuais ou em grupo, listas de exercícios, seminários, atividades práticas e/ou outras estratégias definidas pelo docente, de acordo com a natureza da disciplina.

A composição da nota final (NF) será estabelecida a partir da média ponderada ou aritmética das avaliações aplicadas ao longo do período letivo, conforme critérios previamente apresentados.

\begin{center}
$\mathrm{NF} = \dfrac{A_1 + A_2 + \cdots + A_n}{n}$
\end{center}

ou, quando aplicável,

\begin{center}
$\mathrm{NF} = \sum_{i=1}^{n} w_i \cdot A_i$
\end{center}

em que $A_i$ representa as avaliações realizadas e $w_i$ seus respectivos pesos.

Será considerado aprovado o discente que obtiver frequência mínima de 75\% da carga horária da disciplina e nota final igual ou superior ao mínimo estabelecido pelas normas institucionais vigentes.

Os critérios de avaliação, instrumentos utilizados, pesos atribuídos e cronograma das atividades avaliativas serão apresentados no início do período letivo, podendo ser ajustados, quando necessário, por motivos pedagógicos, mediante comunicação prévia aos discentes.

As atividades avaliativas previstas no plano de ensino, tais como trabalhos individuais ou em grupo, seminários, listas de exercícios, estudos de caso, apresentações e demais instrumentos de acompanhamento da aprendizagem, poderão ser realizadas ao longo do período letivo, conforme cronograma estabelecido pelo docente.

Entretanto, as avaliações na modalidade de prova escrita deverão limitar-se a, no máximo, duas aplicações durante o período acadêmico, devendo ocorrer nas semanas de provas definidas no calendário acadêmico, em conformidade com deliberação aprovada pelo Colegiado da Faculdade de Engenharia Elétrica.%\newpage
% Referências bibliográficas-------------------------------------------
% Programação Semestral (descomente o bloco abaixo para habilitar)-------
%
%  \begin{snugshade}
%  \section*{Programação Semestral} %
%  \end{snugshade}
%
%	\begin{tabular}{ |c|c|p{13cm}| }
%	  \hline
%	  \textbf{Aula}& \textbf{Data}  &  \multicolumn{1}{|c|}{\textbf{Conteúdo}} \\
%	  \hline
%	   01 & 01/01 & Apresentação da disciplina e diretrizes do curso \\
%	  \hline
%	   02 & 01/01 &           \\
%	  \hline
%	   03 & 01/01 &           \\
%	  \hline
%	   04 & 01/01 &           \\
%	  \hline
%	   05 & 01/01 &           \\
%	  \hline
%	   06 & 01/01 &           \\
%	  \hline
%	   07 & 01/01 & Prova Bimestral \\
%	  \hline
%	   08 & 01/01 & Prova Substitutiva \\
%	  \hline
%	   09 & 01/01 &           \\
%	  \hline
%	   10 & 01/01 &           \\
%	  \hline
%	   11 & 01/01 &           \\
%	  \hline
%	   12 & 01/01 &           \\
%	  \hline
%	   13 & 01/01 & Encerramento do conteúdo e revisão de tópicos da disciplina \\
%	  \hline
%	   14 & 01/01 & Prova Final \\
%	  \hline
%	   15 & 01/01 & Prova de Recuperação \\
%      \hline
%
%	\end{tabular}
%\vspace{3mm}
% Referências bibliográficas-------------------------------------------
\vspace{10pt}
\nocite{*}
\nocite{*}

\printbibliography[keyword=basica, title={Bibliografia Básica}]

\printbibliography[keyword=complementar, title={Bibliografia Complementar}]
%\newpage
% Anexos (descomente para habilitar)------------------------------------
%\begin{snugshade}
%\section*{Anexo I} %
%\end{snugshade}
%
%\begin{figure}[h]
%\includegraphics[width=15cm]{Figuras/Universe.png}
%%\centering\caption{There's a map of concepts above.}
%\end{figure}
\end{document}
